% ============================================================
% sec/02trabajos_relacionados.tex
%
% SECCIÓN CENTRAL DE LA ETAPA 1.
%
% DEBE CONTENER:
%   - Eje A: el problema o dominio
%   - Eje B: las técnicas del track
%   - Eje C: el dataset o entorno
%   - Tabla comparativa OBLIGATORIA (7 columnas exactas)
%   - Apartado propio de cierre con las brechas
%
% MÍNIMOS:
%   - 7 artículos arbitrados por integrante
%   - Al menos un survey por integrante
%   - Mayoría de los últimos 5 años
%   - Sin preprints de arXiv salvo justificación aprobada
% ============================================================
\section{Trabajos Relacionados}
\label{sec:relacionados}

\pc{Un párrafo con la metodología de revisión: bases consultadas,
cadenas de búsqueda, criterios de inclusión y exclusión, ventana
temporal, y cuántos trabajos se retuvieron de cuántos revisados.
Esto convierte la revisión en algo auditable en lugar de una lista
de lecturas.}

% ------------------------------------------------------------
% EJE A
% Responde: ¿Cómo se ha abordado este problema en la literatura?
%           ¿Qué enfoques dominan?
%           ¿Qué se considera resuelto y qué sigue abierto?
% ------------------------------------------------------------
\subsection{Eje A: el problema y el dominio}
\label{subsec:ejea}

\pc{3--5 párrafos sintetizando al menos 7 trabajos arbitrados.
AGRUPAR POR FAMILIA DE ENFOQUE, no un párrafo por artículo. Un párrafo
por artículo es un catálogo de lecturas, no un estado del arte.}

\noindent\textbf{Qué enfoques dominan.} \pc{Respuesta directa.}

\noindent\textbf{Qué se considera resuelto.} \pc{Respuesta directa.}

\noindent\textbf{Qué sigue abierto.} \pc{Respuesta directa. Este
párrafo alimenta el apartado de brechas.}

% ------------------------------------------------------------
% EJE B
% Responde: ¿Cuál es el estado actual de los métodos que usaremos?
%           ¿Fortalezas, límites y supuestos bajo los que funcionan?
% AQUÍ SE JUSTIFICA EL TRACK CON LITERATURA, NO CON OPINIÓN.
% ------------------------------------------------------------
\subsection{Eje B: las técnicas del track}
\label{subsec:ejeb}

\pc{3--5 párrafos sintetizando al menos 7 trabajos arbitrados sobre la
familia de técnicas del track elegido.}

\noindent\textbf{Fortalezas.} \pc{Qué resuelven bien estos métodos.}

\noindent\textbf{Límites.} \pc{Dónde fallan.}

\noindent\textbf{Supuestos bajo los que funcionan.} \pc{Qué tiene que
cumplirse en los datos para que el método sea aplicable.}

\noindent\textbf{Justificación técnica del track.} \pc{Por qué este
track y no otro, sostenido con las referencias anteriores. Declarar
explícitamente qué familia alternativa se descartó y con qué argumento
de literatura. Si este párrafo se puede escribir sin citar nada, está
mal escrito.}

% ------------------------------------------------------------
% EJE C
% Responde: ¿Qué trabajos han usado el dataset?
%           ¿Qué métricas reportan y bajo qué protocolo?
%           ¿Cuál es el mejor resultado publicado?
%           ¿Qué problemas conocidos tiene?
% ------------------------------------------------------------
\subsection{Eje C: el dataset}
\label{subsec:ejec}

\pc{3--5 párrafos. Si el dataset propio carece de literatura arbitrada,
DECLARARLO AQUÍ de forma explícita y aplicar la sustitución que permite
el enunciado: revisar datasets análogos del mismo dominio y justificar
por qué son comparables o por qué no lo son.}

\noindent\textbf{Métricas y protocolos reportados.} \pc{Qué mide la
literatura sobre datos análogos y bajo qué partición.}

\noindent\textbf{Problemas conocidos.} \pc{Sesgos, fugas de información,
limitaciones de tamaño o de representatividad. Este apartado es el que
justifica varias decisiones de la Sección~\ref{sec:analisistecnico}.}

% ------------------------------------------------------------
% TABLA COMPARATIVA — OBLIGATORIA
% Columnas exactas: referencia, año, técnica, dataset o entorno,
% métrica principal, resultado reportado, limitación identificada.
% ------------------------------------------------------------
\subsection{Análisis comparativo}
\label{subsec:comparativo}

\safeinput{tab/tabla_comparativa}

\pc{2--3 párrafos LEYENDO la tabla, no repitiéndola: qué tendencia se
ve, qué comparaciones son injustas y por qué, qué limitación se repite
en varios trabajos. Si el texto solo enumera las filas, la tabla no
está cumpliendo su función.}

% ------------------------------------------------------------
% BRECHAS — APARTADO PROPIO DE CIERRE, OBLIGATORIO
% Nombrar carencias concretas, decir cuáles se abordan y cuáles no.
% ------------------------------------------------------------
\subsection{Brechas identificadas}
\label{subsec:brechas}

\textbf{Carencias encontradas en la literatura.}
\begin{enumerate}
  \item \pc{Brecha 1 --- concreta y verificable contra las referencias}
  \item \pc{Brecha 2}
  \item \pc{Brecha 3}
  \item \pc{Brecha 4}
  \item \pc{Brecha 5}
\end{enumerate}

\noindent\textbf{Brechas que este trabajo aborda.} \pc{Cuáles de las
anteriores, y cómo.}

\noindent\textbf{Brechas fuera del alcance.} \pc{Cuáles quedan sin
cubrir y por qué exceden el alcance de este proyecto. Reconocer una
limitación con lucidez es evidencia de dominio, no lo contrario.}